\newpage

\section*{Capítulo 1. Planteamiento del proyecto}
\addcontentsline{toc}{section}{Capítulo 1. Planteamiento del proyecto}
\setcounter{section}{1}
\setcounter{subsection}{0}

\subsection{Problemática}

El mantenimiento de los bosques de la Ciudad de México en los últimos años ha presentado deficiencias notables
debido a la escasez de especialistas y a la falta de tecnologías para su monitoreo. Estos ecosistemas han enfrentado
amenazas constantes por la presencia de plagas, entre ellas \textbf{los insectos descortezadores que en 2022 en promedio
afectaron al 30\% de los árboles de varios bosques de la Ciudad}, que afectan la salud de los árboles y pueden
provocar pérdidas significativas en el entorno natural. Esto representa un problema, ya que, de acuerdo con la
Organización de las Naciones Unidas para la Alimentación y la Agricultura (FAO), los bosques urbanos no solo
contribuyen al equilibrio ecológico, sino que también impactan directamente en la calidad de vida de la población al
mejorar la calidad del aire, reducir la temperatura y ofrecer espacios de recreación \cite{forestal2025bosques}.


En este contexto, los bosques constituyen un elemento fundamental de la biodiversidad en la Ciudad de México. La
presencia de estos insectos representa un riesgo grave, ya que, al alimentarse de la corteza, interrumpen el flujo de
nutrientes, debilitando a los árboles y, en muchos casos, provocando su muerte. Las consecuencias de estas
afectaciones trascienden lo ambiental, alcanzando también dimensiones culturales y recreativas, dado que los bosques
reciben a miles de visitantes diariamente. Además, cada uno de ellos es un espacio emblemático con un profundo valor
histórico y social, ya que albergan museos, áreas deportivas, monumentos y sitios de gran relevancia para la identidad
del país. Su deterioro, por tanto, no solo afecta al equilibrio ecológico, sino también al patrimonio cultural y a la
calidad de los servicios ambientales que proveen.


La detección temprana de este tipo de plagas es esencial para prevenir daños irreversibles y reducir los costos
asociados a la recuperación forestal. Sin embargo, los métodos tradicionales dependen de la inspección visual realizada
por especialistas, lo cual requiere tiempo, experiencia y recursos humanos que no siempre están disponibles. Este
enfoque limita la capacidad de respuesta ante brotes, favoreciendo la propagación de las plagas. A ello se suma que la
extensión de algunos bosques dificulta la supervisión constante de todas sus áreas, lo que aumenta la probabilidad de
que las plagas se desarrollen sin ser identificadas a tiempo.




Ante este panorama, el desarrollo de un sistema móvil que administre y monitoree la plaga en arboles afectados. Esta
herramienta permitirá no solo administrar la información de los árboles afectados, sino también dar un seguimiento
preciso a la evolución de la plaga, todo esto haciendo uso de visión por computadora para que permita identificar
automáticamente plagas de insectos descortezadores siendo una solución más adecuada para apoyar a los especialistas
en el área. Se trata de una oportunidad relevante para innovar y reforzar las estrategias de conservación, buscando el
mantenimiento y recuperación de espacios naturales esenciales en la Ciudad de México. La aplicación de tecnologías
basadas en inteligencia artificial y procesamiento de imágenes podría apoyar el monitoreo continuo, optimizar la toma
de decisiones y generar bases de datos que apoyen a comprender mejor la dinámica de las plagas, posibilitando
estrategias de control más efectivas y sostenibles.



Además, la incorporación de este sistema podría apoyar como modelo para otros espacios verdes del país,
fortaleciendo la gestión ambiental urbana y la preservación de áreas naturales que funcionan como pulmones urbanos
frente a los desafíos del cambio climático y la urbanización acelerada.



